% \documentclass[acmsmall]{acmart}

\documentclass{article}
\usepackage{amssymb}
\usepackage{amsmath}
\usepackage{amsthm}
\usepackage[margin=1.0in]{geometry}
\linespread{1.2}

\usepackage{tikz-cd}
\usepackage{mathpartir}
\usepackage{stmaryrd}
\usepackage{parskip}
\usepackage{xcolor}
\usepackage{array}
\usepackage{comment}

\usepackage{url}
\usepackage{hyperref}

%%%%%%%%%%%%%%%%%%%%%%%%
% theorem environments
%%%%%%%%%%%%%%%%%%%%%%%%
\newtheorem{theorem}{Theorem}[section]
\newtheorem{lemma}{Lemma}[section]
\newtheorem{nonnum-theorem}{Theorem}[section]
\newtheorem{corollary}{Corollary}[section]
\newtheorem{question}{Question}

\theoremstyle{definition}
\newtheorem{definition}{Definition}[section]

\theoremstyle{remark}
\newtheorem{remark}{Remark}
%%%%%%%%%%%%%%%%%%%%%%%%%%%%%%%%%%%%%%%%%%%%%%%%%%%

\newcommand{\To}{\Rightarrow}
\newcommand{\inl}{\mathsf{inl}}
\newcommand{\inr}{\mathsf{inr}}
\newcommand{\alt}{\mathrel{\bf \,\mid\,}}

\newcommand{\ob}{\text{Ob}}


\newcommand{\extlc}{\text{Ext-}\lambda}
\newcommand{\extlcm}{\text{Ext-}\lambda^{-\text{trans}}}
\newcommand{\extlcmm}{\text{Ext-}\lambda^{-\text{trans}-\text{cast}}}
\newcommand{\extlcprime}{\text{Ext-}\lambda'}
\newcommand{\intlc}{\text{Int-}\lambda}
\newcommand{\intlcbisim}{\text{Int$_\approx$-}\lambda}
\newcommand{\erase}[1]{\lfloor {#1} \rfloor}


\newcommand{\uarrowl}{\mathrel{\rotatebox[origin=c]{-30}{$\leftarrowtail$}}}
\newcommand{\uarrowr}{\mathrel{\rotatebox[origin=c]{60}{$\leftarrowtail$}}}
\newcommand{\darrowl}{\mathrel{\rotatebox[origin=c]{30}{$\twoheadleftarrow$}}}
\newcommand{\darrowr}{\mathrel{\rotatebox[origin=c]{120}{$\twoheadleftarrow$}}}
\newcommand{\vuarrow}{\mathrel{\rotatebox[origin=c]{-90}{$\leftarrowtail$}}}
\newcommand{\vdarrow}{\mathrel{\rotatebox[origin=c]{90}{$\twoheadleftarrow$}}}

% Types, terms, and precision 

\newcommand{\dyn}{?}
\newcommand{\nat}{\text{Nat}}
\newcommand{\bool}{\text{Bool}}
\newcommand{\ra}{\rightharpoonup}
\newcommand{\Ret}[1]{\mathsf{Ret\,}{#1}}
\newcommand{\hole}[1]{\bullet \colon {#1}}
\newcommand{\dyntodyn}{\dyn \ra\, \dyn}


\newcommand{\up}[2]{\langle{#2}\uarrowl{#1}\rangle}
\newcommand{\dn}[2]{\langle{#1}\darrowl{#2}\rangle}

\newcommand{\upc}[1]{\text{up}\,{#1}\,}
\newcommand{\dnc}[1]{\text{dn}\,{#1}\,}


% \newcommand{\ret}{\mathsf{ret}}
\newcommand{\err}{\mho}
\newcommand{\zro}{\textsf{zro}}
\newcommand{\suc}{\textsf{suc}}
\newcommand{\lda}[2]{\lambda {#1} . {#2}}

\newcommand{\injarr}[1]{\textsf{Inj}_\ra ({#1})}
\newcommand{\injnat}[1]{\textsf{Inj}_\text{nat} ({#1})}
\newcommand{\casenat}[4]{\text{Case}_\text{nat} ({#1}) \{ \text{no} \to {#2} \alt \text{nat}({#3}) \to {#4} \}}
\newcommand{\casearr}[4]{\text{Case}_\ra ({#1}) \{ \text{no} \to {#2} \alt \text{fun}({#3}) \to {#4} \}}
\newcommand{\casedyn}[5]{\text{Case}_\Dyn ({#1}) \{ \text{nat}({#2}) \to {#3} \alt \text{fun}({#4}) \to {#5} \}}
%\newcommand{\bind}[3]{\text{var } {#1} = {#2} \text{ in } {#3}}
\newcommand{\matchnat}[4]{\text{match } {#1} \text{ with } \{ \textsf {zero} \Rightarrow {#2} \alt \textsf{ suc } {#3} \Rightarrow {#4} \}}

\newcommand{\refl}{\text{refl}}

\newcommand{\Lift}{\text{Lift}}

%\newcommand{\rel}{\circ\hspace{-4px}-\hspace{-4px}\bullet}
\newcommand{\rel}{\mathrel{\circ\mkern-6mu-\mkern-6mu\bullet}}
% \newcommand{\wand}{\mathrel{-\mkern-6mu*}}
\newcommand{\wand}{\mathrel{\multimap}}
\newcommand{\ltdyn}{\sqsubseteq}
\newcommand{\gtdyn}{\sqsupseteq}
\newcommand{\equidyn}{\mathrel{\gtdyn\ltdyn}}
\newcommand{\gamlt}{\Gamma^\ltdyn}
\newcommand{\deltalt}{\Delta^\ltdyn}
\newcommand{\relcomp}{\odot}

\newcommand{\hasty}[3]{{#1} \vdash {#2} \colon {#3}}
\newcommand{\vhasty}[3]{{#1} \vdash^v {#2} \colon {#3}}
\newcommand{\phasty}[3]{{#1} \vdash^c {#2} \colon {#3}}
\newcommand{\etmprec}[4]{{#1} \vdash {#2} \ltdyn_e {#3} \colon {#4}}
\newcommand{\itmprec}[4]{{#1} \vdash {#2} \ltdyn_i {#3} \colon {#4}}
\newcommand{\etmequidyn}[4]{{#1} \vdash {#2} \equidyn_e {#3} \colon {#4}}
\newcommand{\itmequidyn}[4]{{#1} \vdash {#2} \equidyn_i {#3} \colon {#4}}
\newcommand{\synbisim}{\approx_{\text{syn}}}

\newcommand{\Dwn}{\Downarrow}
\newcommand{\qte}[1]{\text{quote}({#1})}

\newcommand{\elab}[1]{\text{Elab}({#1})}

% Perturbations
\newcommand{\pertp}{\text{Pert}^\text{P}}
\newcommand{\perte}{\text{Pert}^\text{E}}
\newcommand{\pertdyn}[2]{\text{pert-dyn}({#1}, {#2})}
\newcommand{\delaypert}[1]{\text{delay-pert}({#1})}

\newcommand{\pertc}{\text{Pert}_{\text{C}}}
\newcommand{\pertv}{\text{Pert}_{\text{V}}}


% SGDT and Intensional Stuff

\newcommand{\later}{{\vartriangleright}}
\newcommand{\laterhs}{{\later}}
\newcommand{\type}{\texttt{Type}}
\newcommand{\lob}{\text{L\"{o}b}}
\newcommand{\tick}{\mathsf{tick}}
\newcommand{\nxt}{\mathsf{next}}
\newcommand{\fix}{\mathsf{fix}}
\newcommand{\kpa}{\kappa}

% Model-related stuff
\newcommand{\calV}{\mathcal{V}}
\newcommand{\calE}{\mathcal{E}}
\newcommand{\calVk}{\mathcal{V}_k}
\newcommand{\calEk}{\mathcal{E}_k}
\newcommand{\tok}{\mathrel{\mathop{\to}\limits^{\textrm{\footnotesize k}}}}
\newcommand{\timesk}{\mathrel{\mathop{\times}\limits^{\textrm{k}}}}
\newcommand{\Fk}{\mathrel{\mathop{F}\limits^{\textrm{k}}}}
\newcommand{\Uk}{\mathrel{\mathop{U}\limits^{\textrm{k}}}}

\newcommand{\op}[1]{{#1}^{\textrm{op}}}
\newcommand{\calC}{\mathcal{C}}
\newcommand{\Set}{\mathsf{Set}}
\newcommand{\ErrDom}{\mathsf{ErrDom}}
\newcommand{\Yo}{\mathsf{Yo}}
\newcommand{\Hom}{\mathsf{Hom}}
\newcommand{\calS}{\mathcal{S}}
\newcommand{\gfix}{\texttt{gfix}}
\newcommand{\qfix}{\texttt{qfix}}
\newcommand{\calU}{\mathcal{U}}
\newcommand{\laterhat}{\widehat{\later}}
\newcommand{\El}{\mathsf{El}}
\newcommand{\Clock}{\mathsf{Clock}}

\newcommand{\Machine}[1]{\mathsf{Machine}\, {#1}}


% Predomains and EP pairs
\newcommand{\Nat}{\mathsf{Nat}}
\newcommand{\Dyn}{\mathsf{Dyn}}
\newcommand{\ty}[1]{\langle {#1} \rangle}
\newcommand{\li}{L_\mho}
\newcommand{\liclk}[1]{L_\mho [{#1}]}

\newcommand{\ext}[2]{\text{ext}\,{#1}\,{#2}}
\newcommand{\map}[2]{\text{map}\,{#1}\,{#2}}

\newcommand{\ltls}{\ltdyn}
\newcommand{\bisim}{\approx}
\newcommand{\semlt}{\le}
\newcommand{\semltbad}{\lesssim}

%\newcommand{\injarr}{\textsf{Inj}_\to}
%\newcommand{\injnat}{\textsf{Inj}_\mathbb{N}}

\newcommand{\id}{\mathsf{id}}
\newcommand{\ep}{\leadsto}

\newcommand{\emb}[2]{\mathsf{emb}_{#1}({#2})}
\newcommand{\proj}[2]{\mathsf{proj}_{#1}({#2})}

\newcommand{\monto}{\to_m}


% Notation for wait functions
\newcommand{\wre}{w_r^e}
\newcommand{\wle}{w_l^e}
\newcommand{\wrp}{w_r^p}
\newcommand{\wlp}{w_l^p}

\newcommand{\sem}[1]{\llbracket {#1} \rrbracket}
\newcommand{\semgl}[1]{{\sem{#1}}^\text{gl}}


% Denotational model
\newcommand{\errdom}{\textsf{ErrDom}}

\newcommand{\ptb}{\text{ptb}}
\newcommand{\push}{\text{push}}
\newcommand{\pull}{\text{pull}}

\newcommand{\pv}{P^{\mathcal{V}}}
\newcommand{\pe}{P^{\mathcal{E}}}
\newcommand{\ptbv}{\text{ptb}^\mathcal{V}}
\newcommand{\ptbe}{\text{ptb}^\mathcal{E}}

\newcommand{\Ppure}{P}
\newcommand{\Pk}{P^K}
\newcommand{\ptbk}{\text{ptb}^K}


\newcommand{\upf}{\text{up}}
\newcommand{\dnf}{\text{dn}
}
\newcommand{\arr}{\to}
\newcommand{\comp}{\bullet}

\newcommand{\ltsq}[2]{\mathrel{\ltdyn^{#1}_{#2}}}
\newcommand{\ltsqbisim}[2]{\mathrel{{\widetilde{\ltdyn}}^{#1}_{#2}}}
\newcommand{\ltbisim}{\mathrel{\widetilde{\ltdyn}}}
\newcommand{\vf}{\mathcal{V}_f}
\newcommand{\vsq}{\mathcal{V}_{sq}}
\newcommand{\ef}{\mathcal{E}_f}
\newcommand{\esq}{\mathcal{E}_{sq}}

\newcommand{\morbisimid}[1]{\text{Endo}_\bisim{#1}}


\newcommand{\sv}{s_{\mathcal{V}}}
\newcommand{\tv}{t_{\mathcal{V}}}
\newcommand{\rv}{r_{\mathcal{V}}}
\newcommand{\se}{s_{\mathcal{E}}}
\newcommand{\te}{t_{\mathcal{E}}}
\newcommand{\re}{r_{\mathcal{E}}}

\newcommand{\Ff}{F_f}
\newcommand{\Fsq}{F_{sq}}
\newcommand{\Uf}{U_f}
\newcommand{\Usq}{U_{sq}}
\newcommand{\arrf}{\arr_f}
\newcommand{\arrsq}{\arr_{sq}}

\newcommand{\vsim}{\mathcal{V}_{\bisim}}
\newcommand{\esim}{\mathcal{E}_{\bisim}}
\newcommand{\svsim}{s_{\mathcal{V}}^\bisim}
\newcommand{\tvsim}{t_{\mathcal{V}}^\bisim}
\newcommand{\rvsim}{r_{\mathcal{V}}^\bisim}
\newcommand{\sesim}{s_{\mathcal{E}}^\bisim}
\newcommand{\tesim}{t_{\mathcal{E}}^\bisim}
\newcommand{\resim}{r_{\mathcal{E}}^\bisim}



\newcommand{\ve}{\mathcal{V}_e}
\newcommand{\ee}{\mathcal{E}_e}

\newcommand{\vr}{\mathcal{V}_r}
\newcommand{\er}{\mathcal{E}_r}

\newcommand{\relatedin}[3]{{#1}\, {#3}\, {#2}}
\newcommand{\binrel}[1]{\mathbin{#1}}


\newcommand{\da}{\downarrow}

\newcommand{\ret}{\text{ret}}
\newcommand{\bind}[3]{{#1} <- {#2}\, ; \, {#3}}

\newcommand{\piv}{\Pi^\mathcal{V}}
\newcommand{\pie}{\Pi^\mathcal{E}}

\newcommand{\upl}{\textsc{UpL}}
\newcommand{\upr}{\textsc{UpR}}
\newcommand{\dnl}{\textsc{DnL}}
\newcommand{\dnr}{\textsc{DnR}}

\newcommand{\delre}{\delta^{r,e}}
\newcommand{\delle}{\delta^{l,e}}
\newcommand{\delrp}{\delta^{r,p}}
\newcommand{\dellp}{\delta^{l,p}}

\newcommand{\qordeq}{\bisim}

\newcommand{\inat}{\text{Inj}_\mathbb{N}}
\newcommand{\itimes}{\text{Inj}_\times}
\newcommand{\iarr}{\text{Inj}_\to}

\newcommand{\tnat}{\mathsf{nat}}
\newcommand{\ttimes}{\mathsf{times}}
\newcommand{\tfun}{\mathsf{fun}}

\newcommand{\cased}[7]{
    \begin{align*}
    \text{Case}_D ({#1}) &\text{ of } \{ \\
        &\alt \text{nat}({#2}) \to {#3}  \\
        &\alt \text{times}({#4}) \to {#5}  \\
        &\alt \text{fun}({#6}) \to {#7} \}
    \end{align*}
}

\newcommand{\inone}{\mathsf{in}_1}
\newcommand{\intwo}{\mathsf{in}_2}
\newcommand{\inthree}{\mathsf{in}_3}
\newcommand{\infour}{\mathsf{in}_4}
\newcommand{\infive}{\mathsf{in}_5}


\newcommand{\delay}{\text{Delay}}
\newcommand{\ledelay}{\le^{\text{Del}}}
\newcommand{\bisimdelay}{\bisim^{\text{Del}}}
\newcommand{\tnow}{\mathsf{now}}
\newcommand{\tlater}{\mathsf{later}}
\newcommand{\Da}{\Downarrow}


\newcommand{\osum}{\text{OSum}}
\newcommand{\inj}[2]{\text{inj}_{#1}({#2})}
\newcommand{\case}[1]{\text{Case}\,{#1}}
\newcommand{\newcase}[3]{\text{newcase}_{#1}{#2}; {#3}}
\newcommand{\icase}[1]{\text{Inj}_{\case{#1}}}
\newcommand{\tref}{\text{ref}}
\newcommand{\new}{\text{new}}
\newcommand{\set}{\text{set}}
\newcommand{\get}{\text{get}}

\definecolor{mygreen}{RGB}{0, 142, 27}
\definecolor{mypurple}{RGB}{90,74,146}
\definecolor{myblue}{RGB}{38,87,175}

\begin{document}

%% TODO
%% - Surface language versus cast calculus
%%
%% - Where to introduce type precision, casts, etc?
%%   1. In the intro?
%%   2. In a syntax section specific to the gradual lambda calculus?
%%   3. In the relational model section?
%%
%% - Too repetitive regarding pitfalls of New-Licata approach
%%
%% - Downcast for the dynamic type involves a case on the sum type

\title{Thesis Proposal}
\author{Eric Giovannini}
\date{April 11, 2025}

\maketitle

\tableofcontents

\section{Introduction}

My thesis work is in the area of \emph{semantics for gradually typed programming
langugaes}. Roughly speaking, gradually-typed languages allow the programmer to
combine both static and dynamic typing disciplines in the same codebase, and to
\emph{gradually} migrate smoothly from a dynamic to static style over time. My
goal has been to advance the state of the art in the mathematical tools used to
model gradually-typed languages and ensure that the languages satisfy the
intended formal properties.
%

In the remainder of this section, we provide background on gradual typing and
the desired properties of gradually-typed languages. We then discuss prior
approaches to formally proving these criteria, with a view towards motivating
the novel approach developed in my thesis work.

\subsection{Gradual Typing: Background}\label{sec:background-gradual-typing}

In programming language design, there is a tension between \emph{static} typing
and \emph{dynamic} typing disciplines. With static typing, the code is
type-checked at compile time, while in dynamic typing, the type checking is
deferred to run-time. Both approaches have benefits and excel in different
scenarios, with static typing offering compile-time assurance of a program's
type safety and type-based reasoning principles that justify program
optimizations, and dynamic typing allowing for rapid prototyping of a codebase
without committing to fixed type signatures.
%
Most languages choose between static or dynamic typing and as a result,
programmers that initially write their code in a dynamically typed language need
to rewrite some or all of their codebase in a static language if they would like
to receive the benefits of static typing once their codebase has matured.

\emph{Gradually typed languages} \cite{siek-taha06, tobin-hochstadt06} seek to
resolve this tension by allowing for both static and dynamic typing disciplines
to be used in the same codebase. These languages support smooth interoperability
between statically-typed and dynamically-typed styles. As a result, the programmer
can begin with fully dynamically-typed code and \emph{gradually} add explicit typing
information to migrate portions of the codebase to a statically typed style
without needing to rewrite the entire project in a completely different language.

% The idea of one program in a gradually-typed language being ``more dynamic'' or
% ``more static'' than another is made formal by the notion of a \emph{term
% precision ordering}. 

The idea of a program becoming ``more static'' over time as the programmer
performs the migration is made formal via a \emph{term precision ordering} on
programs. Term precision gives us the ability to compare programs based on their
relative amount of static typing information. We think of this as a syntactic
relation on programs where a program $M$ is more static or precise than $N$ if,
roughly speaking, $M$ has more typing annotations than $N$.

One of the fundamental theorems for gradually typed languages is
\emph{graduality}, also known as the \emph{dynamic gradual guarantee},
originally defined by Siek, Vitousek, Cimini, and Boyland
\cite{siek_et_al:LIPIcs:2015:5031, new-ahmed2018}.
%
Informally, graduality says that migrating code from dynamic to static typing
should only allow for the introduction of static or dynamic type errors, and not
otherwise change the behavior of the program.
%
% TODO figure here
%
%
This is a way to formalize the programmer's intuition that increasing precision
corresponds to a generalized form of runtime assertions: there are no observable
behavioral changes up to the point of the first dynamic type error.
%
% Fundamentally, this property comes down to the behavior of \emph{runtime type
% casts}, which implement these generalized runtime assertions.

Another desired property of gradually typed languages is that they should offer
some of the benefits of static typing. While standard type soundness, that
well-typed programs are free from runtime errors, is not compatible with runtime
type errors, it is possible instead to prove that gradually typed languages
validate \emph{type-based reasoning}. For instance, while dynamically typed
$\lambda$ calculi only satisfy $\beta$ equality for their type formers,
statically typed $\lambda$ calculi additionally satisfy type-dependent $\eta$
properties that ensure that functions are determined by their behavior under
application and that pattern matching on data types is safe and exhaustive. A
gradually typed calculus that validates these type-dependent $\eta$ laws then
provides some of the type-based reasoning that dynamic languages lack.



\subsection{Prior Work}\label{sec:background-semantics}

%%%%%%%%%%%%%%%%%%%%%%%%%%%%%%%%%%%%%%%%%%%%%%%%%%%%%%%%%%%%%%%%%%

There are several methods of establishing the desired properties of
gradually-typed languages; we discuss these below.

%%\subsubsection{From Static to Gradual}

Siek et al., in the same work that introduces the dynamic gradual guarantee
\cite{siek_et_al:LIPIcs:2015:5031}, also prove the property for a simple
gradually-typed language via a syntactic simulation argument. This proof method
is closely tied to the syntax of the language, and there is limited opportunity
for reuse of common abstractions across different languages.
% and there is limited opportunity for reuse of common building blocks of the
% proofs across different languages.
%
This lack of reusability is addressed somewhat by Siek and Chen
\cite{siek-chen2021}, who give a framework for proving graduality in the Agda
proof assistant. The framework is parameterized by the specific language for
which one wants to prove graduality, and the proof of graduality proceeds via
the syntactic simulation argument.
%
However, the main lemmas of the work are still tied to the particular
operational syntactic calculus being modeled.

%%%%%%%%%%%%%%%%%%%%%%%%%%%%%%%%%%%%%%%%%%%%%%%%%%%%%%%%%%%%%%%%%%

A second line of work for proving graduality consists of starting with a
statically-typed language, and \emph{deriving} from it a language that satisfies
graduality by construction. The methods of Abstracting Gradual Typing
\cite{garcia-clark-tanter2016} and the formal tools of the Gradualizer
\cite{cimini-siek2016} follow this approach. The main downside to these
approaches lies in their inflexibility: since the process in entirely
mechanical, the language designer must adhere to the predefined framework.  Many
gradually typed languages do not fit into either framework, e.g., Typed Racket
\cite{tobin-hochstadt06, tobin-hochstadt08}, and the semantics produced is not
always the desired one.
%
Furthermore, while these frameworks do prove graduality of the resulting
languages, they do not show the correctness of the equational theory, which is
equally important to sound gradual typing.

%%%%%%%%%%%%%%%%%%%%%%%%%%%%%%%%%%%%%%%%%%%%%%%%%%%%%%%%%%%%%%%%%%

% Another major line of work for proving graduality consists of modeling the
% language via an operational semantics and defining a step-indexed logical
% relation over that semantics \cite{new-ahmed2018, new-licata-ahmed2019,
% new-jamner-ahmed19, new-giovannini-licata-2022}.

Another common approach to proving graduality involves defining a semantic model
for a gradually-typed language and establishing the graduality property in the
model.
%
Providing a semantics for gradual typing is inherently complicated in that it
involves (1) recursion and recursive types through the presence of dynamic
typing, (2) effects in the form of divergence and errors, and (3) relational
models to capture the graduality property. Any such semantics needs to address
each of these points.

% A line of work by New, Licata and Ahmed
% \cite{new-ahmed2018,new-licata18,new-licata-ahmed2019} develops an axiomatic
% account of gradual typing and proves graduality by interpreting type precision
% $c : A \ltdyn A'$ as an \emph{embedding-projection pair}, that is, a pure
% embedding function $e : A \to A'$ and a possibly erroring projection $p : \li
% A' \multimap \li A$ such that $p \circ e = \id$ and $e \circ p \ltdyn \id$.

One major line of work in the semantic approach consists of modeling the
language via an operational semantics and defining a step-indexed logical
relation over that semantics \cite{new-ahmed2018, new-licata-ahmed2019,
new-jamner-ahmed19, new-giovannini-licata-2022}. These techniques have been used
to model a variety of languages combining gradual typing with other interesting
features.
%
Step-indexed techniques are robust enough to scale to essentially arbitrary
self-referential definitions, which means they can model highly recursive
programming language features such as higher-order store and runtime-extensible
dynamic types
\cite{appelmcallester01,ahmed06,neis09,new-jamner-ahmed19}.
% Unfortunately, the proofs involve a complicated combination of induction on
% the syntax and tedious step-index manipulation
However, as is common with operational approaches, the graduality proofs based
on these techniques are highly repetitive and technical. Each development
formulates a new logical relation from first-principles and proves many of the
same tedious lemmas without reusable mathematical abstractions across languages.

% TODO mention that I was one of the authors of the GrEff work
% Additionally, the logical relations constructed to prove graduality in prior
% work \cite{new-ahmed2018,new-licata-ahmed2019,new-giovannini-licata-2022} suffer
% from technical complications of requiring two separate expression relations, one
% that counts steps on the left and the other on the right, and there is no
% analogue of this in our approach. However, using two expression relations allows
% some but not all transitive reasoning of term precision to be recovered. In the
% future we aim to explore if this approach is feasible in guarded semantics.


%%%%%%%%%%%%%%%%%%%%%%%%%%%%%%%%%%%%%%%%%%%%%%%%%%%%%%%%%%%%%%%%%%

In the quest for a more compositional and language-independent approach, New and
Licata \cite{new-licata18} constructed a denotational semantics for gradual
typing using the tools of classical domain theory (e.g., $\omega$-CPOs). This
model was extended by Lennon-Bertrand, Maillard, Tabareau and Tanter
\cite{gradualizing-cic} to prove graduality of a gradual dependently typed
calculus $\textrm{CastCIC}^{\mathcal G}$. This domain-theoretic approach
demonstrated that it is possible to give a compositional semantic framework for
proving graduality. However, an approach based on classical domain theory has
fundamental limitations: domain theory appears incapable of modeling certain
perversely recursive features of programming languages such as dynamic type tag
generation and higher-order references \cite{Birkedal-Stovring-Thamsborg-2009},
which are commonplace in real-world gradually typed systems as well as gradual
calculi.


% the inability to model viciously self-referential structures such
% as higher-order extensible state and similar features such as runtime-extensible
% dynamic types.

% Since these features are quite common in dynamically typed languages, we have
% developed in this work a new denotational framework that can be extended to
% model these type system features.

% A further advantage of the denotational approach is that it easily validates
% equational reasoning, not just graduality, and it is completely independent of
% any particular syntax of gradual typing.

%%%%%%%%%%%%%%%%%%%%%%%%%%%%%%%%%%%%%%%%%%%%%%%%%%%%%%%%%%%%%%%%%%

% Goal: combine the compositional nature of the denotational approach  with the
% expressiveness of the step-indexed approach

The goal of the bulk of my thesis work has been to combine the benefits of both
the denotational and step-indexed approaches --- the compositional nature of the
denotational approach and the expressive power of step-indexing. We therefore
take as a starting point the denotational approach of New and Licata. However,
rather than work with classical domain theory, we instead work with
\emph{guarded type theory}, a \emph{synthetic} form of step-indexing. The
synthetic approach to step-indexing introduces additional primitives to the type
theory that capture the essence of step-indexing axiomatically. It is especially
useful as a setting in which to define denotational semantics of languages
involving recursive types, making it a well-suited candidate for modelling
gradually-typed languages.
%
Unfortunately, the compositional theory of New-Licata breaks down in the guarded
setting.
% when accounting for steps, the casts are no longer embedding-projection pairs,
% and we cannot carry out compositional reasoning on terms when those terms
% differ by an arbitrary unknown number of steps. 
My thesis work provides a novel refinement to their semantics that allows for
some compositional reasoning to be recovered even in the guarded setting.

% TODO thesis statement
The claim put forth in this thesis is that \textbf{the use of synthetic methods
facilitates the construction of compositional semantics for real-world
gradually-typed programming languages.}
%
In the remainder of this document, we substantiate this claim by describing the
two major research projects that make up the core of my thesis work. The second
project forms the majority of my thesis work so far and is more directly
relevant to my claim, and we therefore devote significantly more space to
discussing it. We then discuss the next major step in this research program and
outline a plan for the work to be completed.



% Our goal in this work is to provide an \emph{expressive}, \emph{reusable},
% \emph{compositional} semantic framework for defining such well-behaved semantics
% of gradually typed programs.
% %
% Our approach to achieving this goal is to provide a compositional
% \emph{denotational semantics}, mapping types to a kind of semantic domain, terms
% to functions and relations such as term precision to proofs of semantic
% relations between the denoted functions.
% %
% Since the denotational constructions are all syntax-independent, the
% constructions we provide may be reused for similar languages. Since it is
% compositional, components can be mixed and matched depending on what source
% language features are present.



% A line of work by New, Licata and Ahmed
% \cite{new-ahmed2018,new-licata18,new-licata-ahmed2019} develops an axiomatic
% account of gradual typing and proves graduality by interpreting type precision
% $c : A \ltdyn A'$ as an \emph{embedding-projection pair}, that is, a pure
% embedding function $e : A \to A'$ and a possibly erroring projection $p : \li A'
% \multimap \li A$ such that $p \circ e = \id$ and $e \circ p \ltdyn \id$. At
% first glance, our model of precision looks different: it is a \emph{relation}
% $\sem{c} : A \rel A'$ with a quasi-representability structure that gives us
% something close to $e$ and $p$ in an embedding-projection pair. The relationship
% between these models is that \emph{if} the relation were \emph{truly}
% representable we would have that $e$ and $p$ form a \emph{Galois connection} $e
% \circ p \ltdyn \id$ and $\id \ltdyn p \circ e$. However, we have dropped the
% stronger property of retraction from our analysis in this work. With true
% representability, the relation $c$ is \emph{uniquely determined} by the
% embedding $e$, but since we only have quasi-representability, we need to keep
% the relation around explicitly. So the extra complexity of managing explicit
% relations is a cost of the intensional reasoning that guarded type theory
% introduces. 
% %
% The line of work by New, Licata, and Ahmed involved both denotational methods
% and operational logical relations approaches.

% \subsection{Guarded Type Theory}

% \section{Overview of Work Done So Far}

\section{Gradual Typing for Effect Handlers}

My work on gradual typing began with a paper investigating the combination of
gradual typing with algebraic effect handlers \cite{new-giovannini-licata-2022}.
%
In that paper, we presented a new gradually-typed langauge, called GrEff, with
support for algebraic effects and handlers. We sought to study the question of
retrofitting an effect typing system onto an existing language with effects that
lacks such a system. In GrEff, the code in a program can exist along a spectrum
between completely \emph{fully unchecked} and \emph{fully checked} effects. In
the former, the type system is concerned only with the input and output types of
effect operations. With the latter, the type system additionally tracks which
effects a given part of the program may raise.

The motivation for this work was less about the approach to proving graduality
for GrEff, and more about the design of the language itself. Indeed, the method
by which we proved graduality for GrEff is an extension of the operational
logical relations methods employed in prior work. There were some novelties
related to handling gradual subtyping in the semantics, but these were largely
straightforward.

% Maybe start this paragraph differently
However, while constructing these proofs, the lack of reusable building blocks
afforded by the operational approach became especially apparent. Additionally,
manual step-index manipulation was a pervasive part of the development and
proved to be tedious and error-prone. Overall, the time and effort required to
carry out the proofs was enormous: the appendix in the technical report is over
60 pages long \cite{new-giovannini-licata-2022-extended}!
%
We concluded that while the operational approach was sufficient for proving
graduality, the need to reprove all the basic lemmas and essentially redo the
entire development for each new language presents a significant roadblock to
ensuring the soundness of future gradually-typed languages.

% Show the logical relation?

Of course, the operational approach based on step-indexing is not the only means
of establishing graduality. As discussed above, the denotational model of New
and Licata \cite{new-licata18} is not tied to any particular operational
syntactic calculus. Instead, their semantic framework consists of
language-independent, compositional abstractions.
%
% TODO review this part
%
However, by being based on classical domain theory, the semantics of New and
Licata is limited in what language features it can handle. Specifically, it
cannot scale to model languages with higher-order store or dynamic type-tag
generation. Since these features are common in dynamically-typed programming
langugaes, it makes sense that a good semantic foundation for gradual typing
should include support for these features.
% Since these features are common in dynamically-typed programming
% languages, they are also important to include in real-world gradually-typed
% languages. Thus, it makes sense that a good semantic foundation for gradual
% typing should include support for these features.

% We therefore seek to combine the expressive nature of step-indexing with the
% compositional nature of denotational semantics.

We were therefore motivated to seek an alternative approach that features the
best of both worlds; this is the topic of my subsequent project.

\section{Denotational Semantics of Gradual Typing using Synthetic Guarded Domain Theory}

\subsection{Guarded Type Theory}

% In this work, we seek to combine the expressive nature of step-indexing with the
% compositional nature of denotational semantics.
In this work, we seek to define a model of gradual typing that combines the
expressive nature of step-indexing with the compositional nature of denotational
semantics.
%
A key ingredient in achieving this goal is \emph{guarded type theory}, which is
a non-standard foundations with axioms that capture the reasoning principles of
step-indexing \cite{birkedal-mogelberg-schwinghammer-stovring2011,
kristensen-mogelberg-vezzosi2022}. Guarded type theory is a kind of metalanguage
in which we can define denotational semantics of programming languages involving
general recursion. Guarded type theory differs from classical domain theory in
that it can scale to arbitrarily complex recursive definitions.

The new reasoning principles provided by guarded type theory are as follows.
First, there is an operation $\later$ (``later'') that acts on types. A value of
type $\later X$ intuitively represents a value of type $X$ available one
``computational step'' in the future. There is a function $\nxt$ that delays an
element available now to be available later. Lastly, there is \emph{guarded
recursion}. This principle states that there is a unique solution to any
recursive type equation of the form $X \cong F(\later X)$, that is, any
recursive type equation where the variable $X$ is guarded by $\later$.
% (In the step-indexed approach, this corresponds roughly to the idea of
% defining a predicate by induction over the natural number step index.)

The main fact to remember about guarded type theory is that, as with
step-indexing, it is \emph{intensional}, or step-aware: the denotation of a term
in guarded type theory reflects the computational steps it takes. This is in
contrast to classical domain theory, where the steps taken by a term are not
reflected in its denotation. The intensional nature of guarded type theory makes
defining a semantics easier in some places, while introducing challenges in
others.
% State that this makes the relational semantics difficult?

% Much of the interesting work in this paper involves solving the challenges
% arising from the intensional nature of guarded type theory

\subsection{Result}

% Our first major result is a denotational semantics for a simple gradually-typed
% language defined using guarded type theory. The language we model has natural
% numbers, products, and call-by-value-style functions. Furthermore, recall from
% Section \ref{sec:background-semantics} that any semantics for gradually-typed
% languages must model errors and divergence. Therefore, our semantics is a
% partial function from closed terms of type $\nat$ to either a natural number or
% a runtime error:

Our first major result is a denotational semantics for a simple gradually-typed
language defined using guarded type theory. The language we model has natural
numbers, products, and call-by-value-style functions. We will introduce more
details about the language as they are needed. (We refer the reader to our POPL
2025 paper for the full definitions \cite{gtt-sgdt}.) Recall from Section
\ref{sec:background-semantics} that any semantics for gradually-typed languages
must model errors and divergence. Therefore, our semantics is a partial function
from closed terms of type $\nat$ to either a natural number or a runtime error:
%
\[ -\Downarrow : \{M \,|\, \cdot \vdash M : \nat \} \rightharpoonup \mathbb{N}
\cup \{\mho\} \] 
%
It is a partial function because the program may not terminate. Here, $\mho$ is
notation for a runtime type error. We write $M \Downarrow n$ and $M\Downarrow
\mho$ to mean this semantics is defined as a number or error, and $M\Uparrow$ to
mean the semantics is undefined, representing divergence.

Our semantics satisfies the two key properties discussed in Section
\ref{sec:background-gradual-typing}. First, it validates \textbf{type-based
reasoning}. That is, if two closed terms of type $\nat$ are equivalent according
to the equational theory of the language, then they have the same behavior:
either they both evaluate to the same natural number $n$, they both error, or
they both fail to terminate.
%
Second, our semantics satisfies the \textbf{graduality} property: If $M$ and $N$
are closed terms of type $\nat$ where $M$ is more precise than $N$, then either
$M$ errors, in which case we don't care what $N$ does, or $M$ and $N$ have the
same behavior. These results are stated in detail below.

\begin{theorem}[Type-Based Reasoning]
    Let $\cdot \vdash M, N : \nat$. If $M \equiv N$, then $M$ and $N$ have the
    same behavior, in that:
    \begin{align*}
        M \Downarrow n &\iff N \Downarrow n \\
        M \Downarrow \mho &\iff N \Downarrow \mho \\
        M \Uparrow &\iff N \Uparrow
    \end{align*}
\end{theorem}

\begin{theorem}[Graduality]
    Let $\cdot \vdash M, N : \nat$. If $M \ltdyn N$, then either:
    \begin{enumerate}
        \item $M \Downarrow \mho$.
        \item $M \Uparrow$ and $N \Uparrow$.
        \item $M \Downarrow n$ and $N \Downarrow n$ for some natural number $n$.
    \end{enumerate}
\end{theorem}


% A well-behaved semantics should then satisfy several properties. First, it
% should be \emph{adequate}: natural number constants should evaluate to
% themselves: $n \Downarrow n$. Second, it should validate type based reasoning.
% To formalize type based reasoning, languages typically have an equational theory
% $M \equiv N$ specifying when two terms should be considered equivalent.
% Then we want to verify that the big step semantics respects this equational
% theory: if closed programs $M \equiv N$ are equivalent in the equational theory
% then they have the same semantics, $M \Downarrow n \iff N \Downarrow n$ and
% $M\Uparrow \iff N \Uparrow$ and $M \Downarrow \mho \iff N \Downarrow \mho$.


% In the remainder of this section, we will discuss the main ideas in the
% construction of our semantics. We begin with a more detailed overview of the
% gradually-typed language we model.

In the remainder of this section, we will discuss the main ideas in the
construction of the semantics.

% \subsection{Details of the Cast Calculus}
% Give an overview of the language. Include type and term precision?

\subsection{Modeling Gradually-Typed Programs}

We now dive into more detail about our semantics, with an eye towards
understanding the places where guarded type theory makes defining the semantics
simple, and the places where it introduces challenges. 
%
The first question we need to address is how to model effectful programs. Recall
that the language we are modeling involves errors and computational steps. Given
a type $A$, we define a new type we call the \emph{free error domain} on $A$,
written $\li A$. We define the free error domain on $A$ via guarded recursion as
the solution to the equation
%
\[ \li A \cong A + 1 + \later \li A. \]
%
That is, an element of the free error domain is either a returned value of type
$A$, an error, or an element defined ``later''. The latter can be thought of as
a ``thinking'' or ``computing'' element, such that after one time step has
passed we will be able to access the element. Of course, this element may itself
be a value, an error, or another thinking element. 

% We use the notation for the three constructors as shown on the right of the slide.

We will use the notation $\eta : A \to \li A$, $\mho : \li A$, and $\theta :
\later \li A \to \li A$ to refer to the three injections into the sum.

% TODO diverging element

\subsection{Modeling Dynamic Typing}

Next, we consider how to model dynamic typing. Recall that our language, like
all gradually-typed languages, has a special \emph{dynamic type} written $\dyn$
(pronounced ``dyn''). The type $\dyn$ represents a statically-unknown type. The
language we are modeling has natural numbers, products, and call-by-value-style
functions. Thus, to model the dynamic type, we seek a semantic type $D$
satisfying the following equation. 
%
\[ D \cong \Nat + (D \times D) + (D \to \li D). \]
%
Notice that $D$ occurs on the left of the arrow in the function case; this means
that the equation does not have inductive or conductive solutions. It
\emph{does} have solutions using classical domain theory, but in this work we
are using guarded type theory rather than classical domain theory. 
%
Indeed, consider the related equation where we guard the negative occurrence by
placing a later operator outside the function case:
%
\[ D \cong \Nat + (D \times D) + \later (D \to \li D). \]
%
We can solve above equation using a combination of least fixpoint and guarded
recursion. More details on solving the equation are given in our POPL 2025 paper
\cite{gtt-sgdt}.

Though this use of the later modality gives us a semantics for the dynamic type
essentially for free, it comes at a cost: \textbf{A program must take an
observable step whenever it wants to inspect a value of the dynamic type that
happens to be a function.} This has major implications for the relational
semantics, as we discuss below.


\subsection{Relational Semantics}

% Before moving forward, we must introduce some additional important concepts in
% gradual typing. In addition to term precision (introduced in Section
% \ref{sec:background-gradual-typing}), there is also an analogous ordering on
% types called \emph{type precision} that represents when one type is more static
% or dynamic than another. The dynamic type is the top element in this ordering.

Before moving forward, we discuss some additional important concepts in gradual
typing. Recall the notion of term precision introduced in Section
\ref{sec:background-gradual-typing}. There is an analogous ordering on types
called \emph{type precision} that represents when one type is more static or
dynamic than another. Specifically, we write $A \ltdyn A'$ to represent that
type $A$ is more precise (i.e., more static) than type $A'$. We use an explicit
syntax of \emph{type precision derivations} $c : A \ltdyn A'$; see our POPL 2025
paper for the details \cite{gtt-sgdt}. The dynamic type $\dyn$ is the top
element in the type precision ordering, as it is the \emph{least} precise type.
% There is a dynamic type $\dyn$, pronounced ``dyn''. This represents a
% statically-unknown type, and $\dyn$ is the top element in the type precision
% ordering.

Gradually-typed languages are typically elaborated to a lower-level
representation called a \emph{cast calculus}. This involves inserting explicit
coercion terms called \emph{casts} at the boundaries between static and
dynamically-typed code to ensure that the static types are enforced at runtime.
Specifically, whenever $c : A \ltdyn A'$, there is an \emph{upcast} $\upc c$
from terms $A$ to $A'$, and a \emph{downcast} $\dnc c$ from $A'$ to $A$. Upcasts
always succeed, but downcasts can fail at runtime if the types are incompatible.
%
% Gradually-typed languages also feature explicit coercion terms called
% \emph{casts} which are inserted by the compiler at the boundaries between static
% and dynamically-typed code to ensure that the types are enforced at runtime.
% Specifically, whenever $c : A \ltdyn A'$, there is an \emph{upcast} $\upc c$
% from terms $A$ to $A'$, and a \emph{downcast} $\dnc c$ from $A'$ to $A$. Upcasts
% always succeed, but downcasts can fail at runtime if the types are incompatible.
%
% Recall that gradually-typed languages are elaborated to a cast calculus where
% there are explicit casts at the boundaries between static and
% dynamically-typed code. 
The key point to remember about casts in our semantics is that
\textbf{casts can take steps}. For example, downcasting the dynamic type to a
function type involves a step since the function is only available later.

Prior work has established that the graduality property essentially boils down
to the behavior of the cast terms \cite{new-ahmed2018}, shown in Figure
\ref{fig:cast-axioms}. There are four key axioms that capture the intended
behavior of casts, two for upcasts and two for downcasts. These axioms describe
the intended relationship between a type precision and the corresponding
up/downcasts.
%
The basic idea is that \textbf{upcasts make terms less precise}, while
\textbf{downcasts make terms more precise}.

%% TODO figure
\begin{figure}
  \begin{mathpar}

  \inferrule*[right=UpL]
  {M \ltdyn N : c}
  {\upc c M \ltdyn N : A'}

  \inferrule*[right=UpR]
  {M \ltdyn N : A}
  {M \ltdyn \upc c N : c}

  \inferrule*[right=DnL]
  {M \ltdyn N : A'}
  {\dnc c M \ltdyn N : c}

  \inferrule*[right=DnR]
  {M \ltdyn N : c}
  {M \ltdyn \dnc c N : A}

  \end{mathpar}
  \caption{Cast Axioms (Simplified Form)}
  \label{fig:cast-axioms}
\end{figure}

The axioms shown here are simplified from the original formulation as used in
prior work. However, compositional reasoning allows us to \emph{derive} the more
complicated rules from the ones shown here. The alternative would be to prove
the more complicated ones from scratch in the model as was done in prior
operational approaches such as our Gradual Typing for Effect Handlers paper
\cite{new-giovannini-licata-2022}. The derivation of the original cast axioms
from the simplified ones serves as a specific example of the benefits of a
compositional approach to the semantics of gradual typing.

\subsubsection{A First Attempt at a Relational Model}

We next examine how we might attempt to model type and term precision. We will
base the relational semantics on the prior denotational model of New and Licata
\cite{new-licata18}, as follows:

\begin{itemize}
    \item We model our types as sets with a partial ordering relation, representing
    the term precision ordering on values.
    \item Terms will denote monotone functions into the free error domain.
    \item A type precision relationship $c : A \ltdyn A'$ will be modeled as a
    \emph{representable relation} between the denoted posets. The precise
    definition of representability is not important for this discussion; what is
    important is that it allows us to interpret the four cast axioms discussed
    in the previous section.
    \item Term precision will be modeled as a logical-relations style relationship
    between the denoted monotone functions, that is, related inputs are mapped to
    related outputs.
\end{itemize}

Unfortunately, the New-Licata semantics does not carry over directly to the
guarded setting. The issue lies in the requirement that the relations be representable.
%The issue is connected to the definition of representability of relations.
%
In more detail, the problem is a consequence of the tension between the
requirements of \emph{step-insensitivity} and \emph{transitivity}, which are
both necessary properties for our semantics. A relation is step-insensitive when
it is oblivious to the steps taken by terms. We need our relations to be
step-insensitive because on the one hand casts can take steps, but on the other
hand representability requires that casts be related to the identity function,
which does not step. 
%
For example, recall the $\dnr$ rule:
%
\begin{mathpar}
  \inferrule*[right=DnR]
  {M \ltdyn N : c}
  {M \ltdyn \dnc c N : A}
\end{mathpar}
%
Notice that the downcast may take an observable step beyond $N$, but $M$ remains
the same in the premise and conclusion of the rule. In other words, $M$ may
become ``out of sync'' with the downcast.

In addition to step-insensitivity, we also need transitivity. Transitivity is
necessary for us to be able to compose smaller proofs into larger proofs, e.g.,
deriving the original cast axioms from the simplified ones relies on
transitivity.
%
However, as we prove in the paper via a no-go theorem, \textbf{unrestricted
step-insensitivity is not compatible with transitivity}: it is not possible to
define a relation that has both properties.


\subsubsection{Adjusting the Model: Perturbations}

Transitivity is non-negotiable given our goal of a compositional semantics, so
we choose to give up on (unrestricted) step-insensitivity. Specifically, we work
with a \emph{lock-step ordering relation} that is transitive but \emph{sensitive
to steps}. This breaks the representability condition, which as we saw requires
step-insensitivity. (Casts may step and so are not necessarily in lock-step with
the identity function.) To remedy this situation, we restore a small amount of
step-insensitivity in a controlled manner using special synchronizing terms we
call \emph{perturbations}. These are explicitly inserted delay terms that put
both sides into lock-step. Perturbations mimic the stepping behavior of casts,
but are \emph{weakly bisimilar} to the identity function, that is, equivalent to
the identity function when ignoring steps. Finally, we relax the definition of
representability to allow for casts to be lock-step related to perturbations
rather than only to the identity function. We call this notion
\emph{quasi-representability}.

The updated axioms for casts are as shown in Figure \ref{fig:cast-axioms-adjusted}.
Notice that on the sides opposite the casts, we allow for there to be a
\emph{perturbation}. The perturbation mimics the stepping behavior of the cast,
allowing both sides to remain synchronized.

\begin{figure}
  \begin{center}
    \begin{tabular}{m{14em} m{14em}}
      \begin{mathpar}
        \inferrule*[right=UpL]
        {M \ltdyn N : c}
        {\upc c M \ltdyn {\color{blue}\delta^{r,e}(}N{\color{blue}{)}} : A'}
      \end{mathpar} &
      %
      \begin{mathpar}
        \inferrule*[right=UpR]
        {M \ltdyn N : A}
        {{\color{blue}\delta^{l,e}(}M{\color{blue})} \ltdyn \upc c N : c}
      \end{mathpar} \\
      %
      \begin{mathpar}
        \inferrule*[right=DnL]
        {M \ltdyn N : A'}
        {\dnc c M \ltdyn {\color{blue}\delta^{r,p}(}N{\color{blue})} : c}
      \end{mathpar} &
      %
      \begin{mathpar}
        \inferrule*[right=DnR]
        {M \ltdyn N : c}
        {{\color{blue}\delta^{l,p}(} M {\color{blue})} \ltdyn \dnc c N : A}
      \end{mathpar}
    \end{tabular}
  \end{center}

  \caption{Adjusted Cast Axioms}
  \label{fig:cast-axioms-adjusted}
\end{figure}

\subsection{Summary of the Model}

Having introduced perturbations and quasi-representable relations, we now
summarize the final version of the semantics, with the new features
{\color{blue}highlighted in blue}.

\begin{itemize}
  \item Types denote sets with:
  \begin{enumerate}
    \item A partial ordering relation (terms must be in lock-step).
    \item {\color{blue} A weak bisimilarity relation.}
    \item {\color{blue} A monoid of syntactic perturbations.}
  \end{enumerate}
  \item Terms denote \emph{monotone}, {\color{blue}bisimilarity-preserving}
  functions into the free error domain.
  \item Type precision denotes {\color{blue}quasi-}representable relations that
  {\color{blue}interact with syntactic perturbations}.
  \item Term precision is modeled by {\color{blue}taking the lock-step ordering
  on the denoted monotone functions, and closing it under weak bisimilarity on
  both sides.}
\end{itemize}

The last point deserves elaboration: We define the denotation of term precision
$M \ltdyn N$ to mean that there are functions $f$ and $g$ such that $M \bisim f
\ltdyn g \bisim N$ where $\bisim$ is weak bisimilarity and $\ltdyn$ is the
lock-step ordering.

The important takeaway here is that we have separated term precision into a
combination of weak bisimilarity and lock-step reasoning, using syntactic
perturbations as a means to explicitly synchronize terms where needed, which
allows for compositional reasoning while still accommodating the steps taken by
casts.

% \subsection{Adequacy of the Model}


\subsection{Agda Formalization}

In parallel with developing the theory discussed in this paper, we have
developed a partial formalization of our results in Guarded Cubical Agda
\cite{veltri-vezzosi2020}, which extends the Agda type checker with support for
the primitives of guarded type theory. The formalization is available online
\cite{artifact}. Most of the model constructions have been formalized, with the
exception of some lemmas involving quasi-representable relations. We also did
not formalize the translation from the syntax into our semantic model because it
led to prohibitively poor performance of typechecking in Agda.


% \subsection{Investigating Variations of the Semantics}

\section{Extending to More Advanced Language Features}\label{sec:remaining-work}

Our next major step is to extend our semantics to accommodate languages
with higher-order store and/or dynamic type-tag generation.
%
We plan to begin with higher-order store, and then depending on the difficulty
we may also consider dynamic type-tag generation as part of this thesis.

Higher-order store allows for the storage of values of arbitrary type, including
functions. Concretely, there is a type of \emph{references} of type $A$, denoted
$\tref\, A$. The programmer can allocate a new reference cell of type $A$, and
given a reference cell the programmer can perform get and set operations to read
and write the value in the cell.

For the next step of my thesis work, I will consider the language obtained by
augmenting the simple gradually-typed language we modeled previously with the
features of higher-order store. My goal is to adapt our guarded model of gradual
typing to this language. The main technical change that this will require is the
move from a ``set-based'' semantics to a semantics based in \emph{presheaves}
over a category of worlds. This is needed to account for the fact that shape of
the store may change over the course of the program as the programmer allocates
new reference cells. The world tracks the reference cells that have been
allocated, and a presheaf on worlds is essentially a set that is allowed to vary
based on the current world, i.e., the current shape of the store.
%
The main challenge will be to adapt the constructs of our guarded model to the
setting of presheaves. I plan to draw guidance from the work of Sterling et al.
\cite{DBLP:journals/corr/abs-2210-02169}, who give a denotational semantics for
higher-order store using guarded type theory.

% Much of the work will lie in adapting the relational semantics

\begin{comment}

\begin{figure}
    \begin{mathpar}
      % Syntax
      \begin{array}{rcl}
      \text{Types } A &::=& \nat \alt A \ra A' \alt A \times A' \alt \osum \alt {\color{blue} \case{A}}\\
      \text{Type Precision } c &::=&
        r(A) \alt c c' 
             \alt \iarr \alt \inat \alt \itimes 
             \alt {\color{blue} \inj{V}{V}}
             \alt c \ra c' \alt c \times c'\\
      \text{Values } V &::=& 
        x \alt \upc c V \alt \zro \alt \suc\, V 
          \alt \lda{x}{M} \alt (V,V') \\ 
      \text{Terms } M,N &::=& 
        \err \alt \upc c M \alt \dnc c M 
             \alt \zro \alt \suc\, M \alt \lda{x}{M} \\ 
           &&\alt M\, N \alt (M,N) 
             \alt \textrm{let } (x,y) = M \textrm{ in } N
             \alt {\color{blue} \newcase{A}{x}{M} } \\
      \text{Contexts } \Gamma &::= &\cdot \alt \Gamma, x : A \\
      \text{Ctx Precision } \Delta &::=& \cdot\alt \Delta,x:c
    \end{array}
  
    % Selected typing rules
    \inferrule
    {\Gamma \vdash M : A \and c : A \ltdyn A'}
    {\Gamma \vdash \upc c M : A'}
  
    \inferrule
    {\Gamma \vdash N : A' \and c : A \ltdyn A'}
    {\Gamma \vdash \dnc c N : A}

    \inferrule
    {\Gamma, x : \case A \vdash M : A'}
    {\Gamma \vdash \newcase{A}{x}{M} : A'}
  
    \inferrule{}{\Gamma \vdash \mho : A}
    \end{mathpar}
    
    \caption{Cast Calculus Syntax and Typing Rules}
    \label{fig:gtlc-syntax}
  \end{figure}

  \begin{figure}
    % Type precision rules
    \begin{mathpar}
      \inferrule{}{r(A) : A \ltdyn A}

      \inferrule{c : A \ltdyn A' \and c' : A' \ltdyn A''}
                {cc' : A \ltdyn A''}

      \inferrule{}{\iarr \colon \dyn \ra \dyn \ltdyn \dyn}

      \inferrule{}{\inat \colon \nat \ltdyn \dyn}

      \inferrule{}{\itimes \colon \dyn \times \dyn \ltdyn \dyn}

      \inferrule{}{\color{blue}\icase{A} \colon \case{A} \ltdyn \dyn}

      \inferrule{c_i : A_i \ltdyn A_i' \and c_o : A_o \ltdyn A_o'}
                {c_i \ra c_o : (A_i \ra A_o) \ltdyn (A_i' \ra A_o')}

      \inferrule{c_1 : A_1 \ltdyn A_1' \and c_2 : A_2 \ltdyn A_2'}
                {c_1 \times c_2 : (A_1 \times A_2) \ltdyn (A_1' \times A_2')}
    \end{mathpar}

    % Type precision equivalence
    \begin{mathpar}
       r(A)c \equiv c\and
       c \equiv cr(A')\and
       c(c'c'') \equiv (cc')c''\and
       r(A_i \ra A_o) \equiv r(A_i) \ra r(A_o)\and
       r(A_1\times A_2) \equiv r(A_1) \times r(A_2)\and
       (c_i \ra c_o)(c_i' \ra c_o')\equiv (c_ic_i' \ra c_oc_o') \and
       (c_1\times c_2)(c_1'\times c_2')\equiv (c_1c_1' \times c_2c_2')
    \end{mathpar}
    \caption{Type Precision Derivations and Precision Equivalence}
    \label{fig:gtlc-type-precision}
  \end{figure}

\end{comment}


\section{Plan for Remainder of Thesis Work}\label{sec:timeline}

Below is an approximate timeline I am proposing for completing the remainder of
the work. This is subject to change based on the work that we agree upon.

% Note that this is assuming that the technical work remaining is as described in
% the previous section. If we agree that there is additional technical work to be
% done, the timeline will need to change.

\begin{itemize}
    \item May 2025 - October 2025: complete the majority of the remaining
    technical work.
    \item October 2025 - January 2026: wrap-up technical work and focus on
    writing the dissertation.
    % I plan to be a GSI (a graduate teaching assistant) during the fall semester.
    % Therefore, the time I will have to dedicate to the dissertation will be
    % reduced during this period.
    \item January 2026 - April 2026: complete a first draft of the dissertation;
    begin preparing for defense.
    \item May 2026 or August 2026 - defend thesis and complete edits to
    dissertation.
    \item July 2026 - submit the new technical work to the POPL conference.
\end{itemize}

Note: I plan to be a GSI (a graduate teaching assistant) during the Fall 2025
semester. Therefore, the time I will have to dedicate to the dissertation will
be reduced during this period.


% \bibliographystyle{ACM-Reference-Format}
\bibliographystyle{plain}
\bibliography{../../paper-new/references}
% \bibliography{refs}

\end{document}

% In the study of semantics of gradual typing, a distinction is made between the
% \emph{surface language}, which is the language in which the programmer writes
% code, and the \emph{cast calculus}, which is a lower-level language with
% explicit \emph{runtime type casts} that ensure that the static types in a
% gradually-typed program are enforced at runtime. The surface language is
% \emph{elaborated} to the cast calculus and the semantics is defined over the
% cast calculus.